\chapter*{}


\begin{center}
    % Encadré décoratif
    \fboxrule=2pt
    \fboxsep=40pt
    \fcolorbox{black}{gray!7}{
    \centering
        \begin{minipage}[c][0.6\textheight][c]{0.8\textwidth}  % Hauteur = 80% de la page
            \centering

            {\Huge \textbf{Chapitre 1}}\\[1cm]
            {\LARGE Contexte général du projet }\\[0.7cm]
            \rule{0.7\textwidth}{1.5pt}\\[0.7cm]
            {\large Rapport de Projet de Fin d'études}
        \end{minipage}
    }
\end{center}
\vspace*{\fill}

\addcontentsline{toc}{chapter}{Chapitre 1 --  Contexte général du projet}

\setcounter{chapter}{1}
\setcounter{section}{0}
\clearpage

\section*{Introduction}
\addcontentsline{toc}{section}{Introduction}

Dans un contexte de transformation digitale, les institutions bancaires
cherchent à améliorer l'efficacité de leurs traitements internes tout en
garantissant la conformité réglementaire et la qualité des services. Les
opérations du back-office, notamment le traitement des dossiers KYC et des
chèques, représentent à cet égard un enjeu majeur en raison de leur volume,
de leur complexité et de leur caractère répétitif.

Ce chapitre présente le cadre général du projet. Il introduit d'abord
l'organisme d'accueil, puis le contexte et l'étude de l'existant, avant de
formuler la problématique. Il définit ensuite les objectifs visés et les
besoins du système, et décrit enfin la méthodologie et l'approche de
gestion adoptées pour la réalisation du projet.

\section{Présentation de l'organisme d'accueil}

\textbf{Attijariwafa Bank} est l'un des plus grands groupes bancaires et
financiers en Afrique. Né en 2004 de la fusion entre la Banque Commerciale
du Maroc (BCM) et Wafabank, le groupe occupe aujourd'hui une position
stratégique dans le secteur bancaire marocain et africain grâce à son
expertise, sa capacité d'innovation et son réseau étendu.

Présent dans plusieurs pays d'Afrique, d'Europe et du Moyen-Orient, le
groupe accompagne une clientèle diversifiée composée de particuliers, de
professionnels, d'entreprises ainsi que de grandes institutions publiques
et privées. Cette présence internationale lui confère un rôle important
dans le développement économique et financier de plusieurs marchés.

Les activités du groupe couvrent un large ensemble de services financiers
et bancaires, notamment :

\begin{itemize}
    \item La banque de détail ;
    \item La banque des entreprises et des professionnels ;
    \item La banque d'investissement ;
    \item Le financement et les crédits ;
    \item Les services monétiques et moyens de paiement ;
    \item L'assurance et les services financiers spécialisés ;
    \item La gestion des opérations internationales.
\end{itemize}

Dans un environnement bancaire en constante évolution, Attijariwafa Bank
accorde une importance majeure à la transformation digitale et à
l'innovation technologique. Le groupe investit continuellement dans des
solutions numériques modernes afin d'améliorer l'expérience client, de
renforcer la sécurité des opérations et d'optimiser ses processus internes.
C'est précisément dans cette dynamique de modernisation du back-office que
s'inscrit le présent projet de fin d'études.

\subsection*{Organigramme de l'entreprise}

L'organisation d'Attijariwafa Bank repose sur plusieurs directions et pôles
spécialisés collaborant afin d'assurer le bon fonctionnement des activités
bancaires et financières du groupe. La Direction des Systèmes d'Information
y joue un rôle essentiel : elle pilote la mise en œuvre des projets
technologiques et des solutions digitales visant à améliorer les
performances des différents services de la banque.

Le projet réalisé s'intègre principalement dans cette direction, en
interaction avec les services opérationnels chargés du traitement
documentaire et des opérations bancaires.

\begin{figure}[H]
    \centering
    \includegraphics[width=0.85\textwidth]{images/organigramme_atb.png}
    \caption{Organigramme simplifié d'Attijariwafa Bank \cite{organigramme}}
    \label{fig:organigramme}
\end{figure}

La figure ci-dessus présente une vue simplifiée de l'organisation de la
banque ainsi que les différentes directions impliquées dans les activités
liées au projet.

\section{Contexte du projet}

Au sein d'Attijariwafa Bank, plusieurs opérations du back-office reposent
encore sur le traitement manuel de documents administratifs et financiers.
Le projet s'intéresse plus particulièrement à deux de ces processus : les
procédures KYC (\textit{Know Your Customer}) et le traitement des chèques.

Les procédures \textbf{KYC} constituent une étape réglementaire essentielle :
elles permettent de vérifier l'identité des clients et d'assurer la
conformité imposée aux établissements financiers. Dans ce cadre, les clients
fournissent différents documents justificatifs, principalement la carte
nationale d'identité (CIN) et un justificatif de domicile. Le traitement de
ces documents suppose l'identification des informations utiles, le contrôle
de leur conformité, puis leur saisie dans les systèmes internes de la banque.

Le \textbf{traitement des chèques} représente la seconde activité ciblée. Il
implique la lecture des informations portées sur le chèque (montants,
bénéficiaire, coordonnées bancaires), la vérification des montants et de la
signature, ainsi que la validation de l'opération de paiement.

Avec l'augmentation continue du volume des documents traités quotidiennement,
ces traitements manuels atteignent leurs limites face aux exigences de
rapidité, de précision et de conformité. C'est dans ce contexte que
l'intégration de technologies d'intelligence artificielle et de traitement
automatique des documents apparaît comme une réponse adaptée aux besoins de
la banque.

\section{Étude de l'existant}

L'étude de l'existant permet d'analyser le fonctionnement actuel des
processus du back-office, d'en identifier les limites, puis de comparer les
solutions envisageables afin de justifier l'approche retenue.

\subsection*{Description du processus actuel}

Aujourd'hui, les traitements liés aux documents KYC et aux chèques sont
réalisés selon des étapes majoritairement manuelles.

\subsubsection{Traitement des documents KYC}

Les agents du back-office analysent les documents fournis par le client
(CIN et justificatif de domicile) afin d'en extraire les informations
clés : nom et prénom, adresse, numéro de CIN et dates de validité. Ces
informations sont ensuite saisies manuellement dans les systèmes internes
de la banque, après contrôle de leur conformité.

\subsubsection{Traitement des chèques}

Le traitement d'un chèque nécessite plusieurs contrôles successifs : lecture
des informations portées sur le document, vérification des montants (en
chiffres et en lettres), contrôle des coordonnées bancaires, validation de
la signature et vérification de la conformité globale des données. La
majorité de ces vérifications reste réalisée manuellement par les agents.

\subsection*{Critique de l'existant}

Ce fonctionnement, bien qu'opérationnel, présente plusieurs limites :

\begin{itemize}
    \item \textbf{Temps de traitement élevé :} le traitement manuel ralentit
    fortement les opérations lorsque le volume de dossiers augmente, et
    allonge les délais de réponse aux clients.
    \item \textbf{Risque d'erreurs humaines :} les saisies et vérifications
    manuelles sont sources d'erreurs, d'oublis et d'incohérences qui
    affectent la fiabilité des données enregistrées.
    \item \textbf{Diversité des formats documentaires :} les documents
    existent sous de multiples modèles (les justificatifs de domicile, par
    exemple, varient selon l'organisme émetteur), ce qui complique
    l'identification automatique des informations.
    \item \textbf{Charge opérationnelle importante :} les tâches répétitives
    de saisie et de contrôle mobilisent fortement les ressources humaines au
    détriment des tâches à plus forte valeur ajoutée.
\end{itemize}

\subsection*{Étude des solutions existantes}

Pour automatiser ce type de traitement documentaire, plusieurs catégories de
solutions existent sur le marché :

\begin{itemize}
    \item les \textbf{plateformes commerciales d'\textit{Intelligent Document
    Processing} (IDP)}, telles qu'ABBYY FlexiCapture ou Google Document AI,
    qui offrent des performances élevées mais à un coût de licence
    significatif et avec une dépendance à un fournisseur externe ;
    \item les \textbf{services cloud d'OCR} (Amazon Textract, Azure Form
    Recognizer), performants mais soulevant des questions de souveraineté et
    de confidentialité des données bancaires ;
    \item les \textbf{solutions sur mesure} fondées sur des modèles
    \textit{open source} de vision par ordinateur et de reconnaissance de
    texte, qui offrent une maîtrise complète du traitement et des données,
    au prix d'un effort de développement plus important.
\end{itemize}

Compte tenu de la sensibilité des données bancaires et du besoin de maîtrise
totale du pipeline de traitement, c'est cette dernière approche — une
\textbf{solution sur mesure} — qui a été retenue.

\subsection*{Solution proposée}

La solution proposée consiste à concevoir un système intelligent automatisant
une partie des traitements liés aux documents KYC et aux chèques. Elle
permettra notamment :

\begin{itemize}
    \item la détection automatique des zones d'intérêt dans les documents ;
    \item l'extraction automatique des informations présentes sur les CIN et
    les justificatifs de domicile ;
    \item l'assistance au traitement et à la vérification des chèques ;
    \item la réduction des erreurs humaines ;
    \item l'amélioration de la rapidité et de la précision des traitements.
\end{itemize}

Elle s'appuie pour cela sur l'OCR (\textit{Optical Character Recognition}),
les techniques de vision par ordinateur et les modèles de \textit{Machine
Learning} et de \textit{Deep Learning} dédiés à l'analyse documentaire.

\section{Problématique}

Les limites identifiées dans l'étude de l'existant — temps de traitement
élevé, risque d'erreurs humaines, difficulté à absorber un volume croissant
de dossiers et manque de standardisation des processus — convergent vers une
même question centrale :

\begin{center}
\textit{Comment automatiser le traitement des documents du back-office
bancaire à l'aide de l'intelligence artificielle, tout en garantissant la
précision, la rapidité et la conformité des opérations ?}
\end{center}

\section{Objectifs à atteindre}

L'objectif principal de ce projet est de concevoir et de développer une
solution intelligente automatisant les opérations du back-office bancaire,
en particulier le traitement des documents KYC et des chèques, afin
d'améliorer l'efficacité opérationnelle tout en réduisant les tâches
manuelles et les risques d'erreurs.

Cet objectif se décline en plusieurs objectifs spécifiques :

\begin{itemize}
    \item \textbf{Automatiser l'extraction des données :} extraire
    automatiquement les informations des documents bancaires à l'aide de
    l'OCR et de l'intelligence artificielle.
    \item \textbf{Réduire les erreurs humaines :} limiter les erreurs de
    saisie et de vérification grâce à l'automatisation et à des mécanismes de
    contrôle intégrés.
    \item \textbf{Accélérer le traitement des dossiers :} réduire les délais
    d'exécution des opérations et de réponse aux clients.
    \item \textbf{Détecter les anomalies :} identifier automatiquement les
    incohérences, les écarts ou les informations suspectes dans les documents.
    \item \textbf{Mettre en place un système de scoring :} attribuer un score
    de confiance aux dossiers afin d'évaluer leur conformité et de faciliter
    la prise de décision.
    \item \textbf{Améliorer l'efficacité du back-office :} optimiser les
    processus internes en automatisant les tâches répétitives et en assistant
    le travail des agents.
\end{itemize}

\section{Besoins de l'application}

L'identification des besoins permet de définir les fonctionnalités attendues
ainsi que les contraintes techniques et de qualité que doit respecter
l'application. On distingue les besoins fonctionnels des besoins non
fonctionnels.

\subsection*{Besoins fonctionnels}

Le tableau~\ref{tab:besoins_fonctionnels} récapitule les fonctionnalités que
le système doit assurer.

\begin{table}[H]
\centering
\label{tab:besoins_fonctionnels}
\begin{tabular}{|c|l|p{7.5cm}|}
\hline
\textbf{ID} & \textbf{Besoin} & \textbf{Description} \\
\hline
BF1 & Importation des documents & Charger les documents bancaires sous différents formats (images, PDF). \\
\hline
BF2 & Détection des zones d'intérêt & Identifier automatiquement les régions contenant les informations utiles. \\
\hline
BF3 & Extraction des informations & Extraire les données des documents via l'OCR et l'IA. \\
\hline
BF4 & Vérification et validation & Contrôler la cohérence des informations extraites et leur conformité. \\
\hline
BF5 & Traitement KYC & Automatiser la vérification des documents clients et le traitement des dossiers KYC. \\
\hline
BF6 & Traitement des chèques & Analyser les chèques (montants, signature, données bancaires). \\
\hline
BF7 & Affichage des résultats & Présenter les résultats et les données extraites via une interface claire. \\
\hline
BF8 & Score de confiance & Attribuer un score évaluant la fiabilité et la conformité des dossiers. \\
\hline
BF9 & Génération de rapports & Produire des rapports de traitement et de validation. \\
\hline
\end{tabular}
\caption{Besoins fonctionnels du système}
\end{table}

\subsection*{Besoins non fonctionnels}

Les besoins non fonctionnels définissent les critères de qualité du système :

\begin{itemize}
    \item \textbf{Performance :} traitement rapide des documents pour réduire
    les délais.
    \item \textbf{Fiabilité :} précision élevée de l'extraction et de la
    vérification des données.
    \item \textbf{Sécurité :} protection des données bancaires et clients par
    des mécanismes d'authentification et de contrôle d'accès.
    \item \textbf{Scalabilité :} capacité à traiter un volume croissant de
    documents et à évoluer selon les besoins de la banque.
    \item \textbf{Disponibilité :} continuité de service des traitements.
    \item \textbf{Ergonomie :} interface simple et intuitive pour les agents
    du back-office.
    \item \textbf{Maintenabilité :} architecture modulaire facilitant la
    maintenance et les évolutions futures.
\end{itemize}

\section{Approche de travail}

\subsection{Méthodologie adoptée}

Une méthodologie Agile fondée sur le framework \textbf{Scrum} a été adoptée.
Cette approche a permis d'organiser le travail de manière itérative et
incrémentale tout en favorisant l'adaptation continue aux besoins du projet.

Scrum structure le projet en cycles courts appelés \textbf{sprints}, au cours
desquels un ensemble de fonctionnalités est conçu, développé, testé et
validé. Les tâches ont été planifiées et suivies à l'aide d'un \textit{backlog}
regroupant l'ensemble des fonctionnalités à implémenter ; à chaque sprint,
les fonctionnalités prioritaires étaient sélectionnées puis développées
progressivement, ce qui a facilité le suivi de l'avancement et l'intégration
régulière des retours des encadrants.

\begin{figure}[H]
    \centering
    \includegraphics[width=0.85\textwidth]{images/scrum.png}
    \caption{Cycle de vie de la méthodologie Scrum}
    \label{fig:scrum}
\end{figure}

Le processus Scrum suivi s'articule autour des étapes suivantes :

\begin{itemize}
    \item \textbf{Product Backlog :} liste des fonctionnalités et exigences du système.
    \item \textbf{Sprint Planning :} planification des tâches du sprint.
    \item \textbf{Sprint Backlog :} ensemble des tâches sélectionnées pour le sprint en cours.
    \item \textbf{Développement et implémentation :} réalisation des fonctionnalités prévues.
    \item \textbf{Daily Scrum :} suivi quotidien de l'avancement et des blocages.
    \item \textbf{Sprint Review :} présentation des fonctionnalités développées.
    \item \textbf{Sprint Retrospective :} analyse des résultats et identification des axes d'amélioration.
\end{itemize}

\subsection{Formalismes adoptés}

Afin de modéliser et de documenter le système, plusieurs formalismes UML et
de conception ont été utilisés :

\begin{itemize}
    \item \textbf{Diagrammes de cas d'utilisation} pour identifier les
    interactions entre les acteurs et le système ;
    \item \textbf{Diagrammes de classes} pour représenter la structure
    statique de l'application ;
    \item \textbf{Diagrammes de séquence} pour modéliser les échanges entre
    composantes ;
    \item \textbf{Diagrammes d'architecture} pour décrire l'organisation
    globale du système et les interactions entre les services ;
    \item \textbf{Maquettes fonctionnelles} pour visualiser les interfaces
    avant leur implémentation.
\end{itemize}

\section{Gestion de projet}

Le projet a été découpé en sprints successifs, chacun aboutissant à un
livrable concret. Le tableau~\ref{tab:sprints} présente la planification
adoptée.

\begin{table}[H]
\centering
\label{tab:sprints}
\begin{tabular}{|c|p{6.0cm}|p{5.0cm}|}
\hline
\textbf{Sprint} & \textbf{Objectif} & \textbf{Livrable} \\
\hline
Sprint 0 & Étude du contexte, analyse des besoins et choix technologiques & Cahier des charges, backlog initial \\
\hline
Sprint 1 & Conception de l'architecture et des modèles de données & Diagrammes UML, maquettes \\
\hline
Sprint 2 & Pipeline d'extraction des documents KYC (CIN, justificatifs) & Module KYC fonctionnel \\
\hline
Sprint 3 & Pipeline de traitement des chèques & Module chèque fonctionnel \\
\hline
Sprint 4 & Validation, scoring et interface back-office & Application intégrée \\
\hline
Sprint 5 & Tests, corrections et documentation & Version finale et rapport \\
\hline
\end{tabular}
\caption{Planification des sprints du projet}
\end{table}

Le suivi de l'avancement et la répartition des tâches ont été assurés à
l'aide d'un outil de gestion de projet et d'un diagramme de Gantt, permettant
de visualiser l'enchaînement des sprints et le respect des échéances.

% Si tu disposes d'un diagramme de Gantt, ajoute-le ici :
% \begin{figure}[H]
%     \centering
%     \includegraphics[width=0.95\textwidth]{images/gantt.png}
%     \caption{Diagramme de Gantt du projet}
%     \label{fig:gantt}
% \end{figure}

\clearpage

\section*{Conclusion}
\addcontentsline{toc}{section}{Conclusion}

Ce chapitre a posé le cadre général du projet : il a présenté l'organisme
d'accueil, analysé le contexte et l'existant du back-office bancaire, puis
dégagé la problématique. Il a ensuite défini les objectifs, les besoins
fonctionnels et non fonctionnels du système, ainsi que la méthodologie et la
gestion de projet adoptées. Le chapitre suivant sera consacré à la conception
détaillée de la solution proposée.
